\chapter{Reciprocal Pairs, Anchored Stars, and Loop Holonomy}
\label{ch:obstructions}

When does a reciprocal Gaussian agreement between two agents define a law, and what
replaces it when it does not? A tower of scales in which every level's prior is
supplied by the level above has a top without a level above it, and it is tempting
to close the tower by letting two agents supply each other's priors. The general
form of that question is treated in \Cref{sec:full-graph-closing-tower}; this chapter
answers its Gaussian realization. For two latents joined by two parallel edge copies
traversed in opposite directions, the kernel of the assembled precision is, up to a
transport, the fixed space of the loop holonomy, and the determinant factors through
\(\det(I-\mathrm{Hol}(p))\)
(\Cref{prop:obs-reciprocal-pair-kernel,prop:obs-holonomy-determinant-factorization}).
The flat unanchored fold is therefore singular; a proper anchor or holonomy without
eigenvalue one restores definiteness, and in either case the global normalizer
depends on the links and must be retained
(\Cref{prop:obs-normalizer-link-dependence}). The standard conjugate replacement, an
apex latent with a proper prior and conditionally independent constituents, has a
mean-field schedule that converges geometrically to the exact posterior mean at a rate
equal to the largest squared canonical correlation between apex and constituents
(\Cref{thm:obs-star-fixed-point-contraction}). Finally,
\Cref{prop:obs-no-new-anchor-exit} locates the conflict between a nonflat reciprocal
cycle and the strict no-new-anchor coarse family of \Cref{def:cg-no-new-anchor-pair},
a conflict that the default coarse family of this supplement, which admits newly
generated positive-semidefinite self terms, does not have.

Every statement of the chapter concerns a Gaussian realization. The singular fold is a
statement about one flat unanchored reciprocal potential and not about cyclic models,
which are legitimate laws whenever they are globally normalized
(\Cref{sec:gen-undirected}). Graph links follow \Cref{def:geo-graph-links}; the
reciprocal pair uses two distinct parallel edge copies, so its two link values need not
be inverse to one another, while the reversal of each copy carries the inverse link. The
chapter does not assert that the star schedule converges under delayed, noisy,
asynchronous, or inexact updates (\Cref{open:obs-star-update-robustness}). \status{NOT-CLAIMED}
Balanced cut links admit no continuous extension to all singular cuts
(\Cref{open:cg-cut-compression}).

\section{The flat unanchored reciprocal fold}
\label{sec:obs-cycle}

Let two latents \(y_1,y_2\in\R^K\) be joined by distinct edge copies \(e:2\to1\) and
\(e':1\to2\) with links \(\Theta_e,\Theta_{e'}\in\Gamma_b\) and residual covariances
\(\Sigma_e,\Sigma_{e'}\succ0\), so that \(y_2\) supplies a penalty for \(y_1\) and \(y_1\) one
for \(y_2\). The two residuals assemble the precision
\begin{equation}
\Lambda=\delta_e^\top\Sigma_e^{-1}\delta_e+\delta_{e'}^\top\Sigma_{e'}^{-1}\delta_{e'},
\qquad
\delta_e=\begin{pmatrix}I&-\Theta_e\end{pmatrix},
\qquad
\delta_{e'}=\begin{pmatrix}-\Theta_{e'}&I\end{pmatrix}.
\label{eq:obs-pair-precision}
\end{equation}
Let \(p\) be the closed walk that traverses \(e\) and then \(e'\), based at the second
vertex; its holonomy is \(\mathrm{Hol}(p)=\Theta_{e'}\Theta_e\). The reversals of \(e\) and
\(e'\) do not occur in \(p\).

\propositionheading{Kernel of a reciprocal pair}{prop:obs-reciprocal-pair-kernel}
For all residual covariances \(\Sigma_e,\Sigma_{e'}\succ0\),
\begin{equation}
\ker\Lambda=\bigl\{(\Theta_ew,w):w\in\ker(\mathrm{Hol}(p)-I)\bigr\},
\qquad
\dim\ker\Lambda=\dim\ker(\mathrm{Hol}(p)-I).
\label{eq:obs-pair-kernel}
\end{equation}
In particular the kernel does not depend on the residual covariances. \status{ESTABLISHED}

\noindent\emph{Proof.}
\(\Lambda\) is a sum of two positive semidefinite terms, and for \(\Sigma\succ0\) the form
\(v\mapsto v^\top\delta^\top\Sigma^{-1}\delta v\) vanishes exactly when \(\delta v=0\), so
\(\ker\Lambda=\ker\delta_e\cap\ker\delta_{e'}\). For \(v=(y_1,y_2)\) the two conditions are
\(y_1=\Theta_ey_2\) and \(y_2=\Theta_{e'}y_1\), which together say \(y_1=\Theta_ey_2\) and
\(\mathrm{Hol}(p)y_2=y_2\). The map \(w\mapsto(\Theta_ew,w)\) is injective. \hfill$\square$

\medskip
The precision \eqref{eq:obs-pair-precision} is the connection Laplacian of the two-vertex
multigraph with edge copies \(e\) and \(e'\) and weights \(\Sigma_e^{-1}\) and
\(\Sigma_{e'}^{-1}\), and the proposition is its instance of the kernel-holonomy
isomorphism of \Cref{prop:cg-kernel-holonomy}.

\corollaryheading{A flat parallel-edge fold is singular}{cor:obs-flat-fold-singular}
(i) If the two copies give the same flat endpoint comparison,
\(\Theta_{e'}=\Theta_e^{-1}\), then \(\mathrm{Hol}(p)=I\) and
\begin{equation}
\ker\Lambda=\bigl\{(\Theta_ew,w):w\in\R^K\bigr\},
\qquad
\dim\ker\Lambda=K,
\label{eq:obs-cocycle-kernel}
\end{equation}
so the assembled precision is singular for all residual covariances, and its kernel is
the set of transport-consistent configurations.
(ii) For anchors \(A_1,A_2\succ0\),
\begin{equation}
\Lambda_{\mathrm{anch}}=\Lambda+\operatorname{diag}(A_1,A_2)\succ0,
\label{eq:obs-anchored-cycle}
\end{equation}
so for every information vector \(\mathsf h\) the anchored reciprocal energy defines the
normalized Gaussian law with information parameters \((\mathsf h,\Lambda_{\mathrm{anch}})\)
and finite log normalizer \(\mathsf A(\mathsf h,\Lambda_{\mathrm{anch}})\). The same holds for
Gaussian edge energies on any finite multigraph with a positive-definite anchor at every
vertex. \status{ESTABLISHED}

\noindent\emph{Proof.}
(i) is \Cref{prop:obs-reciprocal-pair-kernel} with \(\ker(I-I)=\R^K\). (ii) A sum of
Gaussian edge energies is a positive semidefinite form, and adding a positive-definite
block-diagonal form makes it positive definite; \Cref{prop:gauss-log-normalizer-moments}
gives the normalizer. \hfill$\square$

\medskip
The flat fold charges nothing on the configurations in which the two latents agree under
transport, which is its intended agreement direction, so it supplies no coercivity
there; and an energy whose precision has a \(K\)-dimensional kernel defines no normalized
Lebesgue density, so the evidence identity of \Cref{ch:elbo} has no joint to apply to.
Definiteness returns in two independent ways. A declared anchor works under flat
holonomy. Alternatively, by \eqref{eq:obs-pair-kernel} the unanchored pair is definite
exactly when \(1\notin\operatorname{spec}(\mathrm{Hol}(p))\). That requires
\(\mathrm{Hol}(p)\neq I\), so the two copies cannot have the coboundary form of
\Cref{prop:geo-trivializing-criterion}, which pointwise frame comparisons at one context
always have; at least one copy must carry a link declared independently of the frames in
the sense of \Cref{def:geo-graph-links}. This route leaves the fixed-\(K\) flat interaction
family: a nonflat cycle inside a cluster leaves a self residue
(\Cref{lem:gauss-internal-edge}), and a cut between clusters needs the rectangular
endpoint maps of \Cref{thm:cg-rectangular-endpoint-closure} or the anchored presentations
of \Cref{lem:cg-cut-completed-square}.

\section{Anchors and the global normalizer}
\label{sec:obs-selfprior}

An anchor restores definiteness but not the freedom to treat the normalizer as a
constant. At zero information vector the log normalizer \eqref{eq:gauss-lognorm} of a
precision \(\Lambda\succ0\) on \(\R^n\) reduces to
\begin{equation}
\mathsf A(0,\Lambda)=-\tfrac12\log\det\Lambda+\tfrac n2\log2\pi,
\label{eq:obs-log-normalizer}
\end{equation}
so every dependence of \(\Lambda\) on the links passes into the normalizer. In
\eqref{eq:obs-pair-precision} the links sit inside \(\delta_e\) and \(\delta_{e'}\), and their
contribution to the determinant does not cancel.

\propositionheading{The global normalizer depends on the model graph link}{prop:obs-normalizer-link-dependence}
Take \(K=1\), unit residual variances, flat link values \(\Theta_e=\mathsf g\) and
\(\Theta_{e'}=\mathsf g^{-1}\) with \(\mathsf g>0\), and the anchor \(a_0I\) with
\(a_0>0\). Then
\begin{equation}
\det\bigl(\Lambda+a_0I\bigr)=a_0^2+a_0\bigl(\mathsf g+\mathsf g^{-1}\bigr)^2,
\label{eq:obs-det-witness}
\end{equation}
a nonconstant function of \(\mathsf g\), uniquely minimized on \(\GL^+(1)=(0,\infty)\) at
\(\mathsf g=1\), with value \(a_0^2+4a_0\). At zero information vector the log normalizer
\(\mathsf A(\mathsf g)=-\tfrac12\log\det(\Lambda+a_0I)+\log2\pi\) satisfies
\begin{equation}
\mathsf A'(1)=0,
\qquad
\mathsf A''(1)=-\frac{4}{a_0+4}<0.
\label{eq:obs-normalizer-curvature}
\end{equation}
\status{ESTABLISHED}

\noindent\emph{Proof.}
With \(\delta_e=(1,-\mathsf g)\) and \(\delta_{e'}=(-\mathsf g^{-1},1)\), the two rank-one
contributions sum to
\[
\Lambda=\begin{pmatrix}1+\mathsf g^{-2}&-(\mathsf g+\mathsf g^{-1})\\-(\mathsf g+\mathsf g^{-1})&1+\mathsf g^{2}\end{pmatrix},
\]
whose determinant is \((1+\mathsf g^{-2})(1+\mathsf g^2)-(\mathsf g+\mathsf g^{-1})^2=0\), as
\Cref{cor:obs-flat-fold-singular} requires. Adding \(a_0I\) gives
\(a_0^2+a_0(2+\mathsf g^2+\mathsf g^{-2})\), and \(2+\mathsf g^2+\mathsf g^{-2}=(\mathsf g+\mathsf g^{-1})^2\).
Since \(\mathsf g+\mathsf g^{-1}\geq2\) with equality only at \(\mathsf g=1\), the minimizer is
unique. Writing \(D(\mathsf g)\) for \eqref{eq:obs-det-witness}, \(D'(1)=0\) and
\(D''(1)=8a_0\), so \(\mathsf A''(1)=-\tfrac12D''(1)/D(1)=-4/(a_0+4)\). \hfill$\square$

\medskip
The identity link is thus a strict maximum of the isolated log-normalizer contribution.
The quadratic and information-vector terms also vary with \(\mathsf g\), so no direction
follows for the complete objective; what the witness shows is that the global term
\(-\tfrac12\log\det(\Lambda+a_0I)\) cannot be absorbed into per-agent constants or dropped
from a normalized joint. The variation is a change of model and not of coordinates: it
varies the declared links while holding \(a_0I\), the residual variances, and the Lebesgue
coordinates fixed, whereas under a reframing all of these transform together and
\Cref{prop:gen-product-evidence-invariance} leaves the law and the evidence invariant.

\section{The anchored Gaussian star}
\label{sec:obs-star}

The standard conjugate replacement of the fold is an apex latent inside the joint. Let
\(y_0\in\R^K\) carry the proper prior \(\mathcal N(0,A_0^{-1})\), \(A_0\succ0\), and let \(N\)
constituents be conditionally independent given it,
\begin{equation}
y_0\sim\mathcal N(0,A_0^{-1}),
\qquad
y_i\given y_0\sim\mathcal N(\Theta_iy_0,\Sigma_i),
\qquad
i=1,\dots,N,
\label{eq:obs-star}
\end{equation}
with \(\Theta_i\in\Gamma_b\) and \(\Sigma_i\succ0\). The scale graph is a tree, so no cycle
forms. Write \(\Lambda_\star\) for the joint precision of \((y_0,y_1,\dots,y_N)\). Its
constituent blocks are \(\Sigma_i^{-1}\), its apex-constituent blocks are
\(-\Theta_i^\top\Sigma_i^{-1}\), and its apex block is
\begin{equation}
\Lambda_{00}=A_0+\sum_{i=1}^N\Theta_i^\top\Sigma_i^{-1}\Theta_i .
\label{eq:obs-precision-addition}
\end{equation}
This is the conditional precision of \(y_0\) given the constituents in the linear-Gaussian
form of Bayes' theorem \citep[Sec.~2.3.3]{Bishop2006}, and it is the precision of the
optimal mean-field factor of the apex, which for a Gaussian joint is the corresponding
diagonal block of the joint precision \citep[Sec.~10.1.2]{Bishop2006}. The constituents
enter the apex through the single statistic \(\sum_i\Theta_i^\top\Sigma_i^{-1}y_i\), whose
conditional covariance given \(y_0\) is \(\Lambda_{00}-A_0\). The star is one conjugate
hierarchical Gaussian model among many; the computations below use only its tree form
and its proper root.

\propositionheading{The declared prior supplies exactly the missing coercivity}{prop:obs-anchor-coercivity}
For all invertible \(\Theta_i\) and all \(\Sigma_i\succ0\), \(\Lambda_\star\succ0\). On every
transport-consistent configuration \(w=(v,\Theta_1v,\dots,\Theta_Nv)\),
\begin{equation}
w^\top\Lambda_\star w=v^\top A_0v
\label{eq:obs-coercivity}
\end{equation}
for every \(v\in\R^K\). \status{ESTABLISHED}

\noindent\emph{Proof.}
The quadratic form is
\(y_0^\top A_0y_0+\sum_i(y_i-\Theta_iy_0)^\top\Sigma_i^{-1}(y_i-\Theta_iy_0)\), a sum of
positive semidefinite terms. It vanishes only if \(y_0=0\), because \(A_0\succ0\), and
then \(y_i=\Theta_iy_0=0\). On a transport-consistent configuration every residual
vanishes and only \(v^\top A_0v\) remains. \hfill$\square$

\medskip
The contrast with the fold is exact: by \Cref{cor:obs-flat-fold-singular} the flat pair
charges nothing on its transport-consistent directions, and by \eqref{eq:obs-coercivity} the
star charges exactly \(v^\top A_0v\) there.

\theoremheading{Exact Gaussian-star fixed point and contraction}{thm:obs-star-fixed-point-contraction}
Let \(A_0\succeq0\), \(O_i\succeq0\), and let \(r_0,r_1,\dots,r_N\in\R^K\). Consider the
Gaussian target on \((y_0,y_1,\dots,y_N)\) whose log density is, up to a constant,
\[
-\tfrac12y_0^\top A_0y_0
-\tfrac12\sum_i(y_i-\Theta_iy_0)^\top\Sigma_i^{-1}(y_i-\Theta_iy_0)
-\tfrac12\sum_iy_i^\top O_iy_i
+r_0^\top y_0+\sum_ir_i^\top y_i ;
\]
the terms in \(O_i\) and \(r_i\) are Gaussian observations of the constituents, and
\(A_0\succ0\), \(O_i=0\), \(r=0\) recover \eqref{eq:obs-star}. Keep \(\Lambda_{00}\) of
\eqref{eq:obs-precision-addition}, write \(\Lambda_{ii}=\Sigma_i^{-1}+O_i\), and put
\begin{equation}
B_\star=\sum_{i=1}^N\Theta_i^\top\Sigma_i^{-1}\Lambda_{ii}^{-1}\Sigma_i^{-1}\Theta_i,
\qquad
S_\star=\Lambda_{00}-B_\star,
\qquad
M=\Lambda_{00}^{-1}B_\star,
\qquad
\bar r=r_0+\sum_{i=1}^N\Theta_i^\top\Sigma_i^{-1}\Lambda_{ii}^{-1}r_i .
\label{eq:obs-star-contraction-data}
\end{equation}
Assume \(S_\star\succ0\), which holds whenever \(A_0\succ0\) and is equivalent to positive
definiteness of the target precision. Consider the mean-field schedule that updates every
constituent factor from the current apex factor, in parallel or in any order, and then
updates the apex factor. The factor covariances become \(\Lambda_{ii}^{-1}\) and
\(\Lambda_{00}^{-1}\) on their first updates, and the means obey
\begin{align}
\mu_i^{(t+1)}
&=\Lambda_{ii}^{-1}\left(\Sigma_i^{-1}\Theta_i\mu_0^{(t)}+r_i\right),
\label{eq:obs-star-constituent-update}\\
\mu_0^{(t+1)}
&=\Lambda_{00}^{-1}\left(r_0+\sum_i\Theta_i^\top\Sigma_i^{-1}\mu_i^{(t+1)}\right)
=\Lambda_{00}^{-1}\bar r+M\mu_0^{(t)}.
\label{eq:obs-star-apex-update}
\end{align}
There is a unique fixed point,
\begin{equation}
\mu_0^\star=S_\star^{-1}\bar r,
\qquad
\mu_i^\star=\Lambda_{ii}^{-1}\left(\Sigma_i^{-1}\Theta_i\mu_0^\star+r_i\right),
\label{eq:obs-star-fixed-point}
\end{equation}
which is the exact mean of the target, and in the norm
\(\lVert x\rVert_{\Lambda_{00}}=(x^\top\Lambda_{00}x)^{1/2}\),
\begin{equation}
\lVert\mu_0^{(t+1)}-\mu_0^\star\rVert_{\Lambda_{00}}
\leq\varrho_\star
\lVert\mu_0^{(t)}-\mu_0^\star\rVert_{\Lambda_{00}},
\qquad
\varrho_\star=\lambda_{\max}\left(\Lambda_{00}^{-1/2}B_\star\Lambda_{00}^{-1/2}\right)<1 .
\label{eq:obs-star-rate}
\end{equation}
The rate \(\varrho_\star\) is the largest squared canonical correlation between \(y_0\) and
\((y_1,\dots,y_N)\) under the target. The relative entropy of the apex factor to its
fixed-point factor contracts by \(\varrho_\star^2\) per sweep, and the product factor
converges to the product of the fixed-point factors in relative entropy. When every
\(O_i=0\), one has \(B_\star=\Lambda_{00}-A_0\) and \(S_\star=A_0\), so the hypothesis is
\(A_0\succ0\) and the fixed point is
\(\mu_0^\star=A_0^{-1}(r_0+\sum_i\Theta_i^\top r_i)\), \(\mu_i^\star=\Theta_i\mu_0^\star+\Sigma_ir_i\);
if instead \(A_0\succeq0\) is singular and every \(O_i=0\), then \(\varrho_\star=1\) and
neither uniqueness nor contraction holds. \status{ESTABLISHED}

\noindent\emph{Proof.}
The target precision \(\Lambda_\star\) has apex block \(\Lambda_{00}\), constituent blocks
\(\Lambda_{ii}\), apex-constituent blocks \(-\Theta_i^\top\Sigma_i^{-1}\), and no block between
distinct constituents. Its constituent part \(\operatorname{diag}(\Lambda_{ii})\) is positive
definite, and the Schur complement of that part is \(S_\star\), so \(\Lambda_\star\succ0\)
exactly when \(S_\star\succ0\). If \(A_0\succ0\), then
\(S_\star-A_0=\sum_i\Theta_i^\top(\Sigma_i^{-1}-\Sigma_i^{-1}\Lambda_{ii}^{-1}\Sigma_i^{-1})\Theta_i\succeq0\),
because \(\Lambda_{ii}^{-1}\preceq\Sigma_i\). The optimal mean-field factor of one block is
Gaussian with precision the diagonal block of \(\Lambda_\star\) and with the information
vector of the target minus the off-diagonal blocks applied to the current means of the
other factors \citep[Secs.~10.1.1--10.1.2]{Bishop2006}; this gives the covariances,
\eqref{eq:obs-star-constituent-update}, and the first expression in
\eqref{eq:obs-star-apex-update}, and substitution gives the affine recursion. Since
\(I-M=\Lambda_{00}^{-1}S_\star\) is invertible, the recursion has the unique fixed point
\eqref{eq:obs-star-fixed-point}; the two block equations it satisfies are the two block rows
of \(\Lambda_\star\mu=r\), so it is the exact mean.

Because \(\Lambda_{00}\succeq B_\star+S_\star\) with \(S_\star\succ0\) and \(B_\star\succeq0\), the matrix
\(M\) is similar to the symmetric matrix
\(\Lambda_{00}^{-1/2}B_\star\Lambda_{00}^{-1/2}=I-\Lambda_{00}^{-1/2}S_\star\Lambda_{00}^{-1/2}\), whose
eigenvalues lie in \([0,1)\). With \(\xi_t=\mu_0^{(t)}-\mu_0^\star\),
\(\lVert\xi_{t+1}\rVert_{\Lambda_{00}}=\lVert\Lambda_{00}^{-1/2}B_\star\Lambda_{00}^{-1/2}\Lambda_{00}^{1/2}\xi_t\rVert_2
\leq\varrho_\star\lVert\xi_t\rVert_{\Lambda_{00}}\), which is \eqref{eq:obs-star-rate}. Write
\(\Lambda_{0y}\), \(\Lambda_{yy}\), and \(\Sigma_\star=\Lambda_\star^{-1}\) with blocks
\(\Sigma_{00},\Sigma_{0y},\Sigma_{yy}\) for the apex and the stacked constituents. Then
\(B_\star=\Lambda_{0y}\Lambda_{yy}^{-1}\Lambda_{y0}\), the block identities of
\(\Lambda_\star\Sigma_\star=I\) give \(-\Lambda_{00}^{-1}\Lambda_{0y}=\Sigma_{0y}\Sigma_{yy}^{-1}\) and
\(-\Lambda_{yy}^{-1}\Lambda_{y0}=\Sigma_{y0}\Sigma_{00}^{-1}\), and so
\(M=\Sigma_{0y}\Sigma_{yy}^{-1}\Sigma_{y0}\Sigma_{00}^{-1}\), whose eigenvalues are the squared
canonical correlations between \(y_0\) and the constituents.

After the first updates every factor has its fixed-point covariance. The apex factor then
satisfies \(\KL(q_0^{(t)}\Vert q_0^\star)=\tfrac12\lVert\xi_t\rVert_{\Lambda_{00}}^2\), which
contracts by \(\varrho_\star^2\) per sweep. The constituent means updated from
\(\mu_0^{(t)}\) differ from their fixed points by \(\Lambda_{ii}^{-1}\Sigma_i^{-1}\Theta_i\xi_t\),
so the sum of their divergences is \(\tfrac12\xi_t^\top B_\star\xi_t\leq\tfrac12\xi_t^\top\Lambda_{00}\xi_t\),
and the product factor converges in relative entropy. When every \(O_i=0\),
\(\Lambda_{ii}=\Sigma_i^{-1}\) gives \(B_\star=\sum_i\Theta_i^\top\Sigma_i^{-1}\Theta_i\) and
\(S_\star=A_0\), and the displayed fixed point follows. If then \(A_0v=0\) for some
\(v\neq0\), \(Mv=\Lambda_{00}^{-1}(\Lambda_{00}-A_0)v=v\), so \(\varrho_\star=1\), \(I-M\) is singular,
and the error along \(v\) is never reduced. \hfill$\square$

\medskip
The mean recursion is block Gauss--Seidel on \(\Lambda_\star\mu=r\). Convergence of the point
Gauss--Seidel iteration for a symmetric positive definite system is classical
\citep[Thm.~11.2.3]{GolubVanLoan2013}, and the same iteration matrix governs the rate of the
deterministic-scan Gibbs sampler for a Gaussian target \citep{RobertsSahu1997}. What the
theorem supplies for the star is the closed form of \(\varrho_\star\), its identification
with a canonical correlation, the extension to observed constituents, and the convergence
of the product factor in relative entropy.

\openproblemheading{Robustness of the star schedule}{open:obs-star-update-robustness}
\Cref{thm:obs-star-fixed-point-contraction} covers the exact block schedule in which every
constituent update reads the current apex factor and the apex update reads the updated
constituents. Whether the fixed point is reached, and at what rate, when the updates are
delayed, noisy, asynchronous, or inexact is open. A convergence theorem for a declared
class of perturbed schedules, or a counterexample within such a class, would settle it.

\section{Loop holonomy and the reciprocal fold}
\label{sec:obs-flatness}

By \Cref{prop:obs-reciprocal-pair-kernel} the obstruction of \Cref{sec:obs-cycle} is
measured by one number, \(\dim\ker(\mathrm{Hol}(p)-I)\), and
\Cref{cor:obs-flat-fold-singular} evaluates it under flat holonomy, where it is largest.
The kernel is produced by the flat loop product, not by the cycle.

\corollaryheading{Nontrivial holonomy shrinks the reciprocal-pair kernel}{cor:obs-holonomy-kernel-shrinkage}
Under the hypotheses of \Cref{prop:obs-reciprocal-pair-kernel},
\begin{equation}
\begin{aligned}
\dim\ker\Lambda=K
&\iff\mathrm{Hol}(p)=I,\\
\mathrm{Hol}(p)\neq I
&\implies\dim\ker\Lambda\leq K-1,\\
\ker\Lambda=\{0\}
&\iff1\notin\operatorname{spec}(\mathrm{Hol}(p)).
\end{aligned}
\label{eq:obs-holonomy-kernel}
\end{equation}
Under the specialization \(\Gamma_b=\GL^+(K)\), the holonomies at which the last condition
fails form the intersection of \(\GL^+(K)\) with the zero set of the nonzero polynomial
\(\det(\mathrm{Hol}-I)\), a relatively closed subset with dense complement and Lebesgue measure
zero, so the flat point \(\mathrm{Hol}(p)=I\) is the most degenerate member of the family and a
generic holonomy gives no kernel. If instead \(\Gamma_b=\operatorname{SO}(K)\) with \(K\) odd,
every holonomy has eigenvalue one and the unanchored pair is singular for every choice of
links. \status{ESTABLISHED}

\noindent\emph{Proof.}
By \eqref{eq:obs-pair-kernel}, \(\dim\ker\Lambda=\dim\ker(\mathrm{Hol}(p)-I)\), which is \(K\)
exactly when \(\mathrm{Hol}(p)-I=0\), at most \(K-1\) otherwise, and zero exactly when one is not
an eigenvalue. The polynomial \(\det(\mathrm{Hol}-I)\) takes the value \(1\) at
\(\mathrm{Hol}=2I\), and the zero set of a nonzero polynomial is closed, nowhere dense, and
Lebesgue null. An element of \(\operatorname{SO}(K)\) has determinant one and eigenvalues of
modulus one in conjugate pairs; for odd \(K\) the number of real eigenvalues is odd and their
product is one, so at least one of them equals one. \hfill$\square$

\propositionheading{Graph-link dependence of the determinant factors through holonomy}{prop:obs-holonomy-determinant-factorization}
For the reciprocal pair \eqref{eq:obs-pair-precision},
\begin{equation}
\det\Lambda=\frac{\bigl(\det(I-\mathrm{Hol}(p))\bigr)^{2}}{\det\Sigma_e\det\Sigma_{e'}} .
\label{eq:obs-holonomy-det}
\end{equation}
The links enter the determinant only through the loop holonomy, \(\Lambda\) is singular exactly
when \(\det(I-\mathrm{Hol}(p))=0\), and where \(\Lambda\succ0\) the link-dependent part of the
zero-information log normalizer \eqref{eq:obs-log-normalizer} is
\begin{equation}
-\tfrac12\log\det\Lambda
=-\log\bigl\lvert\det(I-\mathrm{Hol}(p))\bigr\rvert
+\tfrac12\log\det\Sigma_e
+\tfrac12\log\det\Sigma_{e'} .
\label{eq:obs-holonomy-normalizer}
\end{equation}
The singular locus \(\{\det(I-\mathrm{Hol}(p))=0\}\) strictly contains the flat point: for
\(K=2\), \(\Theta_e=I\), \(\Theta_{e'}=\operatorname{diag}(1,2)\), and
\(\Sigma_e=\Sigma_{e'}=I\), one has \(\mathrm{Hol}(p)\neq I\) and \(\det\Lambda=0\).
\status{ESTABLISHED}

\noindent\emph{Proof.}
Stack the residual maps as
\(\delta=\bigl(\begin{smallmatrix}I&-\Theta_e\\-\Theta_{e'}&I\end{smallmatrix}\bigr)\), so that
\(\Lambda=\delta^\top\operatorname{diag}(\Sigma_e,\Sigma_{e'})^{-1}\delta\) by
\eqref{eq:obs-pair-precision}. Hence
\(\det\Lambda=(\det\delta)^2/(\det\Sigma_e\det\Sigma_{e'})\), and the Schur formula gives
\(\det\delta=\det(I-\Theta_{e'}\Theta_e)\). Taking logarithms gives
\eqref{eq:obs-holonomy-normalizer}. In the witness \(\mathrm{Hol}(p)=\operatorname{diag}(1,2)\) and
\(\det(I-\mathrm{Hol}(p))=0\). \hfill$\square$

\medskip
Holonomy without eigenvalue one removes exactly one defect of this architecture, the
kernel of the unanchored pair: the reciprocal energy with its global normalizer then
defines a Gaussian law, and the transport-consistent zero modes collapse to the origin.
The statement concerns the two-residual precision \eqref{eq:obs-pair-precision} and no other
cyclic potential. Holonomy does not turn the reciprocal energy into a pair of locally
normalized directed kernels; it defines an undirected Gaussian potential whose global
partition function is part of the model, and by \eqref{eq:obs-holonomy-normalizer} that term
is a nonconstant function of the links that diverges on the singular locus, where the
unanchored law ceases to exist. Holonomy also does not remove declarations. A closed walk
with nontrivial holonomy requires at least one link declared independently of the frames,
because links of the coboundary form of \Cref{prop:geo-trivializing-criterion} telescope to
the identity around every closed walk, and pointwise frame comparisons at one context
(\Cref{sec:geo-regime-one}) have that form; replacing a declared prior by a declared link
relocates the declaration. Both \Cref{prop:obs-normalizer-link-dependence} and
\eqref{eq:obs-holonomy-normalizer} vary declared links while holding the other coordinate
parameters fixed, so neither is a statement about motion along a gauge orbit.

The strict no-new-anchor pair subfamily of \Cref{def:cg-no-new-anchor-pair}, in which a
blocking step creates no self term beyond the summed self terms and the internal residues,
conflicts with an unanchored reciprocal pair. Its exact fixed-\(K\) aggregation needs
clusters whose links can all be removed by a choice of frames, and exact closure of each
cut by one bare coarse link, while the unanchored pair needs definiteness:
\begin{equation}
\begin{aligned}
\text{strict flat coarse-graining needs}\quad
&\text{trivial holonomy on every closed walk in its domain},\\
\text{the unanchored reciprocal pair needs}\quad
&1\notin\operatorname{spec}(\mathrm{Hol}(p)).
\end{aligned}
\label{eq:obs-tension}
\end{equation}

\propositionheading{A nonflat cycle leaves the strict flat coarse-graining}{prop:obs-no-new-anchor-exit}
Let a finite sequence of nested hard partitions of the agents be given in which every
cluster at every level is trivializing, that is, the links inside it take the coboundary
form of \Cref{prop:geo-trivializing-criterion}. Suppose that the coarsest level is either one
cluster or two clusters whose cut, with positive-definite weights, is reproduced exactly by
one coarse link without a new self term, as in the subfamily of
\Cref{def:cg-no-new-anchor-pair}. Then every closed walk of the interaction graph has trivial
holonomy. In particular, no such sequence exists for a population that contains an
unanchored reciprocal pair with \(1\notin\operatorname{spec}(\mathrm{Hol}(p))\).
\status{ESTABLISHED}

\noindent\emph{Proof.}
If the coarsest level is one cluster, that cluster is the whole population and is
trivializing by hypothesis. If it is two clusters \(I\) and \(J\), choose frames in which every
link inside \(I\) and inside \(J\) is the identity. The cut links become twists \(\Theta_e\in\Gamma_b\),
and exact reproduction by one coarse link without a new self term is \(\Delta=0\) in
\eqref{eq:cg-cut-excess}, which for positive-definite weights holds exactly when all twists
are equal. Re-expressing every frame of \(J\) by this common twist makes every link of the
population the identity, so the whole population is trivializing and every closed walk has
trivial holonomy. For the reciprocal pair this gives \(\mathrm{Hol}(p)=I\), which has eigenvalue
one. \hfill$\square$

\medskip
Within the strict subfamily, a nonflat reciprocal cycle can therefore stay outside the flat
coarse-graining only by stopping the sequence before it reaches the coarsest levels of the
proposition. The default coarse family of this supplement admits newly generated
positive-semidefinite self terms, and there the conflict dissolves. For a cut between two
trivialized clusters with weights \(W_e\succ0\), twists \(\Theta_e\), and \(\Lambda_{II}\), \(\Lambda_{IJ}\),
\(\bar\Theta=\Lambda_{II}^{-1}\Lambda_{IJ}\), and \(\Delta\) as in \eqref{eq:cg-cut-excess}, \Cref{lem:cg-cut-completed-square} writes the cut
energy as
\begin{equation}
\sum_e(z_I-\Theta_ez_J)^{\top}W_e(z_I-\Theta_ez_J)
=(z_I-\bar\Theta z_J)^{\top}\Lambda_{II}(z_I-\bar\Theta z_J)+z_J^{\top}\Delta z_J ,
\label{eq:obs-cut-completed-square}
\end{equation}
one link \(\bar\Theta\) of weight \(\Lambda_{II}\) together with the self term \(\Delta\) at \(J\), and an exact
anchored one-link presentation exists if and only if \(\det\Lambda_{IJ}\neq0\), with \(\det\Lambda_{IJ}>0\) for
links in \(\GL^+(K)\). The admitted self term is exactly the anchor that removes the
right-hand demand of \eqref{eq:obs-tension}; the balanced presentation of
\Cref{prop:cg-balanced-cut-presentation} selects one of the presentations, no presentation at the
level of laws is balanced once the cut has excess (\Cref{prop:cg-no-balanced-law-presentation}), and
no continuous extension of the balanced links covers all singular and rectangular cuts
(\Cref{open:cg-cut-compression}).
